\chapter{Cohomologie de la retenue hélicoïdale : évaluation des cocycles et dimension minimale de la mémoire cumulative}
\label{chap:elevacion-cohomologica-dimension-rev2}
\index{dimension!rang cohomologique}
\index{retenue!intégrale discrète}
\index{mémoire!relèvement dimensionnel}

Une coordonnée périodique enregistre la position visible ; la retenue enregistre
le nombre de fois où la trajectoire a franchi sa frontière. La première revient,
la seconde s'accumule. Cette différence transforme la périodicité en un
relèvement : deux instants peuvent partager une phase visible et appartenir à des états
distincts parce qu'ils conservent des nombres de tours différents.

Dans l'alphabet nonadique, pour \(n\in\mathbb Z_{\geq1}\), définissons
\[
 r(n)=1+((n-1)\bmod9),
 \qquad
 q(n)=\left\lfloor\frac{n-1}{9}\right\rfloor,
\]
et le relèvement hélicoïdal
\begin{equation}
 H_9(n)=(r(n),q(n)).
 \label{eq:elevacion-helicoidal-nonadica-rev2}
\end{equation}
Alors
\[
 H_9(n+9)=H_9(n)+(0,1):
\]
la phase revient et la mémoire change. Si \(c_k\in\mathbb Z\) est la retenue
du pas \(k\to k+1\), son intégrale discrète
\begin{equation}
 Q_N=Q_0+\sum_{k=0}^{N-1}c_k
 \label{eq:integral-discreta-carry-rev2}
\end{equation}
constitue la coordonnée cumulative du relèvement. Son indépendance est
algébrique ; une réalisation métrique ultérieure peut la représenter par
une direction transversale ou orthogonale.

\section{Cocycles indépendants et nombre minimal de coordonnées de mémoire}

Soit \(G\) le graphe orienté des états visibles et soient
\(c_1,\ldots,c_d\in Z^1(G;\mathbb Z)\) des cocycles entiers. Pour une route
\(\gamma\), définissons
\[
 Q_j(\gamma)=\sum_{e\in\gamma}c_j(e),
 \qquad
 \widetilde\gamma=
 \bigl(q_{\rm vis}(\gamma),Q_1(\gamma),\ldots,Q_d(\gamma)\bigr).
\]
Sur une route ouverte, cette évaluation retient le représentant du cocycle
et les extrémités. En effet, avec la convention
\((\mathrm d\phi)(e)=\phi(t(e))-\phi(s(e))\),
\begin{equation}
 Q_j^{\,c_j+\mathrm d\phi}(\gamma)
 =Q_j^{\,c_j}(\gamma)+\phi(t\gamma)-\phi(s\gamma).
 \label{eq:cambio-representante-cociclo-ruta-rev2}
\end{equation}
Par conséquent, pour une route ouverte, on fixe le représentant et ses extrémités,
et la mémoire conserve ces données. Si \(\gamma\) est un cycle, le terme de
frontière s'annule et \(Q_j(\gamma)\) dépend uniquement de la classe
\([c_j]\in H^1(G;\mathbb Z)\).

\begin{theorem}[Critère cohomologique de rang dimensionnel]
\label{thm:rango-dimensional-cohomologico-rev2}
Supposons que les classes de \(c_1,\ldots,c_d\) engendrent un sous-module libre de
rang \(d\) dans \(H^1(G;\mathbb Z)\). Soit
\[
 Q:Z_1(G;\mathbb Z)\longrightarrow\mathbb Z^d
\]
leur évaluation sur les cycles. Si une famille de \(m\) compteurs additifs s'
obtient par un homomorphisme \(A:\mathbb Z^d\to\mathbb Z^m\), de sorte
que
\[
 \widehat Q:=A\circ Q:Z_1(G;\mathbb Z)\longrightarrow\mathbb Z^m,
\]
et permet de reconstruire les évaluations, c'est-à-dire qu'il existe
\(B:\mathbb Z^m\to\mathbb Z^d\) avec
\(Q=B\circ\widehat Q\), alors
\[
 \boxed{m\geq d}.
\]
Par conséquent, le relèvement nécessite au moins \(d\) coordonnées de mémoire
indépendantes en plus de la coordonnée visible.
\end{theorem}

\begin{proof}
L'accouplement
\(H^1(G;\mathbb Z)\times H_1(G;\mathbb Z)\to\mathbb Z\) et l'indépendance
de \([c_1],\ldots,[c_d]\) impliquent
\(\operatorname{rk}\operatorname{im}Q=d\). Comme
\(Q=B\circ\widehat Q\),
\[
 d=\operatorname{rk}\operatorname{im}Q
 \leq\operatorname{rk}\operatorname{im}\widehat Q\leq m.
\]
Aucun homomorphisme vers moins de \(d\) compteurs entiers n'admet d'
inverse additif sur le sous-module engendré par les évaluations \(Q_j\).
La condition plus forte \(B\circ A=I_{\mathbb Z^d}\) est suffisante, mais la
borne ne nécessite que la reconstruction de \(Q\) sur les évaluations
réalisées par les cycles.
\end{proof}

La dimension généalogique quantifie ici les degrés de liberté nécessaires pour
retrouver le transport. Le théorème n'identifie pas encore ce rang à une
dimension spatiale physique : il fournit le domaine minimal sur lequel une
réalisation métrique peut agir sans effacer les cocycles indépendants.

\section{Tour nonadique et conservation projective de l'unité}

Soit
\[
 \mathcal C_9^{(d)}=(\mathbb Z/9\mathbb Z)^d.
\]
Pour une marque \(\star\in\mathbb Z/9\mathbb Z\), nous définissons l'inclusion
\[
 h_{d,\star}:\mathcal C_9^{(d)}\longrightarrow
 \mathcal C_9^{(d+1)},
 \qquad h_{d,\star}(x)=(x,\star),
\]
et la projection
\[
 \pi_{d+1,d}(x_1,\ldots,x_{d+1})=(x_1,\ldots,x_d).
\]

\begin{proposition}[Conservation par relèvement dimensionnel]
\label{prop:conservacion-elevacion-dimensional-rev2}
On a
\[
 \pi_{d+1,d}\circ h_{d,\star}=\operatorname{id}_{\mathcal C_9^{(d)}}.
\]
Par conséquent, la cellule de dimension \(d\) réapparaît exactement comme une
section coordonnée de la cellule de dimension \(d+1\) : une tranche affine, et un
sous-groupe lorsque \(\star=0\). Le terme \emph{face de codimension un} est
réservé à une réalisation cubique supplémentaire expressément déclarée. Pour les mesures
uniformes normalisées \(\mu_d\),
\[
 (\pi_{d+1,d})_*\mu_{d+1}=\mu_d.
\]
\end{proposition}

\begin{proof}
La première identité résulte de l'élimination de la dernière coordonnée tout juste
ajoutée. Pour un point \(x\in\mathcal C_9^{(d)}\), la fibre de
\(\pi_{d+1,d}\) contient neuf points, chacun de masse
\(9^{-(d+1)}\) ; sa masse totale est \(9^{-d}=\mu_d(\{x\})\).
\end{proof}

La première égalité exprime la conservation structurelle ; la seconde, la
conservation normalisée. Relever ne remplace pas la couche précédente : cela l'inclut
comme section et distribue la nouvelle coordonnée sur ses fibres. En dimension
trois, \(|\mathcal C_9^{(3)}|=729\) ; sa complétion décimale et la
décomposition disjointe
\[
 1000=729+243+27+1
\]
demeurent dans la résidence de la complétion cubique, où sont distingués
l'intérieur, les faces, les arêtes et le sommet ajouté.

\section{De la coordonnée cumulative à la métrique}

Le relèvement des cocycles et le choix d'une métrique sont des opérations
successives. Le premier construit des coordonnées indépendantes de mémoire ; le
second déclare une forme quadratique qui mesure leurs variations. Dans APP, le
bloc central
\[
 B_c=\begin{pmatrix}7&2\\2&7\end{pmatrix}=9P_++5P_-
\]
produit la normalisation de référence
\[
 M_c:=\frac{B_c}{\sqrt{45}}\in SO^+(1,1),
\]
démontrée dans le \cref{prop:app-trit-espectral-markov-lorentz-rev8}.
Dans l'écran cubique, la réalisation physique sélectionne la
signature lorentzienne correspondante et le lecteur temporel de TPK oriente sa
dynamique. La chaîne causale est
\[
 \begin{aligned}
 \text{retour avec retenue}
 &\longrightarrow\text{cocycle cumulatif}
 \longrightarrow\text{rang de mémoire}\\
 &\longrightarrow\text{forme quadratique}
 \longrightarrow\text{métrique réalisée}.
 \end{aligned}
\]

\paragraph{Séparation entre le rang cohomologique interne et sa réalisation spatio--temporelle.}
L'hélice nonadique et l'intégrale de la retenue fournissent les prémisses du
critère cohomologique de rang et de sa borne minimale. La tour de cellules et la
conservation de la mesure se déduisent de la projection finie. L'identification
du rang généalogique à une dimension spatio-temporelle requiert le foncteur
métrique déclaré en mécanique dimensionnelle.
